\section{Назначение и сценарии работы}

\subsection{Назначение продукта}

Основной продукт --- мобильный сканер для складских операций в Л'Этуаль. Он помогает сотруднику открыть подготовленное задание, быстро считать штрихкод товара, проверить данные и подтвердить результат. Справочник товаров, ячейки и задания загружаются в приложение из подготовленных файлов импорта.

\subsection{Пользователи и доступы}

\begin{tabularx}{\textwidth}{L{45mm}Y}
\toprule
\textbf{Роль} & \textbf{Доступ и действия} \\
\midrule
Складской сотрудник & Открывает локально загруженные задания на приёмку или размещение, сканирует товары, проверяет данные и подтверждает выполнение. Видит сведения, необходимые для выполнения текущей операции. \\
\bottomrule
\end{tabularx}

\subsection{Причина разработки}

Используемые устройства работают с устаревшим ПО, которое сложно сопровождать и расширять. Требуется Java-приложение со сканером, которое легко сопровождать и расширять.

\subsection{Границы прототипа и ограничения}\label{scope:limitations}

Цель прототипа --- показать мобильный сканер штрихкодов и понятный порядок действий сотрудника при складских операциях. Приложение считывает код товара или коробки, показывает относящееся к операции задание, дает сотруднику проверить данные и фиксирует подтвержденный результат. Приоритетом являются скорость и надежность сканирования, простой интерфейс для ТСД и независимость бизнес-сценариев от конкретной реализации сканера.

В состав прототипа \textbf{входят}:

\begin{itemize}
  \item экран сканирования штрихкода с обработкой успешного считывания, ошибки и отказа в доступе к камере;
  \item единый интерфейс модуля сканирования, который в дальнейшем позволит заменить камерный сканер на Zebra DataWedge, SDK ТСД или другую реализацию без изменения логики операций.
  \item проверка отсканированного кода в рамках открытого задания и явное подтверждение результата сотрудником;
  \item импорт подготовленных файлов со справочником товаров, ячейками и заданиями;
  \item хранение каталога, заданий, результатов операций и отчётов о приёмке только на устройстве;
  \item сохранение фотографии подписанной накладной в отчёте о приёмке для последующего просмотра документа;
\end{itemize}

В состав прототипа \textbf{не входят}:

\begin{itemize}
  \item интеграция с 1С, ERP, WMS, API поставщиков, корпоративными сервисами или любыми другими внешними учетными системами;
  \item автоматическая синхронизация справочника товаров, остатков, заданий и результатов операций с сервером;
  \item разработка собственной полноценной системы учета товаров, управления остатками, заказами, поставками или складскими ячейками;
  \item серверная часть бизнес-операций, многопользовательская работа, управление ролями, авторизация и централизованное администрирование;
  \item юридически значимый электронный документооборот, электронная подпись и хранение оригиналов документов;
\end{itemize}


\subsection{Общий цикл работы с JavaТСД}

\begin{enumerate}
  \item Сотрудник выбирает подготовленное задание в приложении.
  \item Открывает задание и видит тип операции, зону или ячейку, товар и требуемое действие.
  \item Сканирует EAN-13 товара, при необходимости указывает фактическое количество или подтверждает ячейку.
  \item Подтверждает выполнение; приложение фиксирует результат и сообщает об ошибках или расхождениях.
  \item Приложение сохраняет и показывает итог операции на устройстве. Передача результатов во внешнюю учётную систему не выполняется.
\end{enumerate}
